\subsection{Define the Quantity}

The first step in any measurement campaign is to clearly define the quantity or phenomenon to be measured or observed. Once this is established, a method for observing or quantifying that quantity must be designed.

In metrology, this is often referred to as an \emph{indirect measurement}, since in many cases we cannot directly observe the quantity of interest. A simple example is measuring electrical current. There are no methods to directly measure current in a straightforward way; instead, we measure a related, observable quantity. Common approaches include:

\begin{itemize}
\item Using a shunt resistor to measure the voltage drop across it.
\item Employing a Hall-effect sensor to measure the magnetic field generated by the current.
\item Measuring the magnetic flux through a coil induced by the current.
\end{itemize}

In each case, we ultimately measure a voltage that is related to the current. This is considered an indirect measurement, where the desired quantity (current) is inferred from a relationship to a directly observable quantity (voltage, magnetic field, or flux).

Although the previous example is simple, the principle of indirect measurement applies broadly. For instance, one might wish to observe the RF spectrum or the reflection coefficients that arise when a hand is moved near an antenna. In this case, the indirect measurement is the reflection coefficient, while the desired phenomenon is the placement or motion of the hand. By measuring the reflection coefficient, we infer the effect of the hand on the antenna, illustrating how indirect measurements allow us to quantify phenomena that cannot be directly observed.

It is important to note that in exploratory work - such as investigating whether the position of a hand can be inferred from measured reflection coefficients - the principle of indirect measurement still applies. In this case, we must not only measure the reflection coefficients, but also know the exact placement of the hand. This allows us to map a relationship between the observable quantity (reflection coefficients) and the underlying phenomenon (hand position), which is essential for drawing meaningful conclusions.

It is also acknowledged that some experimental phenomena are difficult to quantify directly. For example, the observed effect may not be the exact placement of a hand, but rather a hand gesture or motion. In such cases, it may not be meaningful to assign a precise uncertainty to the degree or type of gesture.

The key principle is to remain aware of these limitations and to clearly describe them in the dataset documentation. Even if the phenomenon cannot be rigorously quantified, explicitly stating the context, assumptions, and limitations ensures that subsequent users of the data understand its scope and reliability.
